\begin{pastexam}[label={px:nccu-113-dsa-02}]{政大資科所 113 學年度 第 2 題}
（單選；不必說明理由）

(4\%) Which of the following sorting algorithms has a best-case time complexity of $O(n)$?
\begin{enumerate}[label=\Alph*.]
\item Merge sort
\item Insertion sort
\item Quick sort
\item Heap sort
\end{enumerate}
\end{pastexam}

\begin{pastsoln}
答案是 B。

\begin{itemize}
\item[A.] 合併排序不看資料長什麼樣子，每次都切兩半、合併，遞迴式永遠是 $T(n) = 2T(n/2) + \Theta(n)$，最好也是
$\Theta(n \lg n)$（資結書第~14.4~節）。 合併時就算一邊先用完，另一邊直接接上，比較次數至少也有 $n/2$ 量級，
加總起來還是 $\Theta(n \lg n)$。
\item[B.] 插入排序在輸入已經排好時，每個數只和前一個比一次就停，共 $n - 1$ 次比較，$\Theta(n)$
（資結書第~1.3、1.8~節）。 對。
\item[C.] 快速排序每一層分割都要把整段掃過一次；最好的情況是每次切成兩半，$\Theta(n \lg n)$
（資結書第~14.5~節）。
\item[D.] 堆積排序要做 $n - 1$ 次取出最大值，鍵互不相同時，最好情況也是 $\Theta(n \lg n)$
（資結書第~14.6~節）。
\end{itemize}
補充：C、D 在「所有鍵都相同」這種特殊輸入下，用三路分割的快速排序或堆積排序可以是 $O(n)$；題目問的是一般
課本上的最好情況，指鍵互不相同的輸入，答案只有插入排序。
\end{pastsoln}
